\GalSourcePageBoundary{085}{L0084}{56}{56}
de la forme
\[
(a_{k_{1},k_{2}},\ a_{k_{1},\phi k_{2}})
\]
et de substitutions de l'ordre $p^{2} - p$, dont la période serait de $p$~termes.
(\emph{Voyez} encore l'endroit cité.)

Les premiers doivent nécessairement, pour que le groupe
jouisse de la propriété voulue, se réduire à la forme
\[
(a_{k_{1},k_{2}},\ a_{k_{1},m k_{2}})
\]
d'après ce qu'on a vu pour les équations de degré~$p$.

Quant aux substitutions dont la période serait de $p$~termes,
comme elles sont conjuguées aux précédentes, nous pouvons supposer
un groupe qui les contienne sans contenir celles-ci: donc
elles devront transformer les substitutions circulaires~\textup{(a)} les unes
dans les autres; donc elles seront aussi linéaires.

Nous sommes donc arrivés à cette conclusion, que le groupe
primitif de permutations de $p^{2}$~lettres doit ne contenir que des
substitutions de la forme~\textup{(A)}.

Maintenant, prenons le groupe total que l'on obtient en opérant
sur l'expression
\[
a_{k_{1},k_{2}}
\]
toutes les substitutions linéaires possibles, et cherchons quels
sont les diviseurs de ce groupe qui peuvent jouir de la propriété
voulue pour la résolubilité des équations.

Quel est d'abord le nombre total des substitutions linéaires?
Premièrement, il est clair que toute transformation de la forme
\[
k_{1},\ k_{2},\quad m_{1}k_{1} + n_{1}k_{2} + \alpha_{1},\ m_{2}k_{1} + n_{2}k_{2} + \alpha_{2}
\]
ne sera pas pour cela une substitution; car il faut, dans une substitution,
qu'à chaque lettre de la première permutation il ne réponde
qu'une seule lettre de la seconde, et réciproquement.

Si donc on prend une lettre quelconque $a_{l_{1},l_{2}}$ de la seconde permutation,
et que l'on remonte à la lettre correspondante dans la
première, on devra trouver une lettre $a_{k_{1},k_{2}}$ où les indices $k_{1}$,~$k_{2}$
seront parfaitement déterminés. Il faut donc que, quels que soient
$l_{1}$~et~$l_{2}$, on ait, par les deux équations
\[
m_{1}k_{1} + n_{1}k_{2} + \alpha_{1} = l_{1}, \quad
m_{2}k_{1} + n_{2}k_{2} + \alpha_{2} = l_{2},
\]
